% Cette fiche d'activité à été développée par
% + Ismaël Gaye <ismael.gaye@ens-rennes.fr>
% + Adil Oubninte <adil.oubninte@ens-rennes.fr>

% Le template a été largement inspiré de celui de
% + Erwan Tanguy-Legac <erwan.tanguy-legac@ens-rennes.fr>
% + Héloïse Troublé <heloise.trouble@ens-rennes.fr>

% Elle est distribuée sous la licence Creative-Commons by-nc-sa 4.0
% (Attribution, Non-Commercial, Share Alike)
% qui peut être trouvée dans
% [le site de Creative Commons](https://creativecommons.org/licenses/by-nc-sa/4.0/)

\documentclass[12pt, a4paper]{article}
\usepackage[french]{babel}
\usepackage[utf8]{inputenc}
\usepackage[top=0.9in, left=0.6in, right=0.7in, bottom=1in]{geometry}
\usepackage{amssymb}
\usepackage{color}
\usepackage{fancyhdr}
\pagestyle{fancy}
\usepackage{longtable}
\usepackage{array}
\usepackage{titling}

\renewcommand{\headrulewidth}{1pt}
\chead{} 
\lhead{Alice déménage - Plan de l'activité}
\rhead{Ismaël Gaye et Adil Oubninte}

\renewcommand{\footrulewidth}{1pt}
\cfoot{\thepage} 
\lfoot{2024}
\rfoot{ENS Rennes}

\fancypagestyle{plain}{
	\fancyhf{} 
	\cfoot{\thepage}
	\renewcommand{\headrulewidth}{0pt}
	\renewcommand{\footrulewidth}{1pt}}

\begin{document}
	
	\title{Alice déménage}
	\author{Ismaël Gaye et Adil Oubninte}
	\date{24 mars 2024}
	\maketitle
	
	\paragraph*{Niveau :} $6^{\grave{e}me}$
	\paragraph*{Durée :} 55 minutes
	\paragraph*{Prérequis :} Aucun
	\paragraph*{Description du jeu :}
	Alice déménage est un jeu se jouant tout seul mais peut aussi se faire en coopération. L'objectif de ce jeu est étant donnés une grille de départ (représentant la taille maximale de la pile que Alice peut faire au dessus de sa voiture) et de ses valises (ici des polygones) de construire la pile de valises de hauteur minimale.
	
	\paragraph*{Objectif pédagogique :}
	Introduire la notion de difficulté des problèmes (et donc la NP-complétude).
	
	\begin{longtable}{c|m{3.9cm}|m{8.4cm}|m{2.4cm}}		\textbf{Durée} & \textbf{Phase} & \textbf{Activités et consignes} & \textbf{Organisation et matériel} \\ \hline\hline
		
		5' &
		
		Introduction
		
		&
		
		On se présente (rapidement cette fois), on explique les règles. On fait un exemple au tableau. On insiste bien sur le fait de trouver un algorithme qu'un ordinateur peut comprendre plutôt que de trouver une solution (Les tables seront positionnées en îlot en amont de la séance)
		
		& 
		
		A l'oral au tableau
		
		\\ \hline
		
		10' & 
		
		Temps libre
		
		&
		
		On distribue les grilles et première enveloppe (qui ne seront pas sur les tables au début de la séance). On laisse les élèves essayer de trouver un moyen de trouver la solution optimale, on s'assure que tous ont bien compris les règles, on répond aux questions (possibilité de gagner du temps en cas de retard à l'installation)
		
		& 
		
		Groupes de 2 (possiblement en îlots de 4 pour échanger), 12 jetons par groupe (+ext pour les plus rapides)
		
		\\ \hline

        5' &
		
		Remise en commun
		
		&
		
		On stoppe l'activité, on reprend l'attention et on ouvre une discussion tous ensemble. Si possible, on interroge des groupes qui ont trouvé une esquise d'algorithme. On fait venir des élèves aux tableaux pour les tester sur l'exemple.
		
		& 
		
		À l'oral, au tableau
		
		\\ \hline
		
		15'
		
		&
		
		Robot
		
		& 
		
		Les élèves continue de jouer mais cette fois l'un d'entre eux doit donner des instructions élémentaire (pour simuler un algorithme). Au besoin on passe dans les groupes pour faire le robot.
		
		& Par groupe
		
		\\  \hline
		
		10'
		
		&
		
		Robot efficace
		
		&
		
		On distribue tout le jeux au élèves et on lui demande d'avoir un robot efficace, qui ne dépend pas du jeux.
		
		& 
		
		Par groupes de 4 (les îlots)
		
		\\ \hline
		
		10'
		
		&
		
		Institutionnalisation
		
		&
		
		On reprend l'attention des élèves et on explique l'insitutionnalisation de manière intéractive. On répond aux potentiels question et on s'assure qu'ils ont compris en leur demandant de ré-expliquer
		
		&
		
		A l'oral
		
	\end{longtable}
	
\newpage
	
\subsection*{Institutionnalisation :}

Le problème d'arranger les affaires d'Alice sur le toit de sa voiture, qui est de largeur fixe est "compliqué": il faut en général beaucoup d'étape pour trouver la meilleure solution (et ça grandit très vite quand le nombre d'affaire augmente)
Un ordinateur aussi mettra beaucoup de temps, car certain groupe l'ont fait, il doit construire la meilleure solution en retournant sur ses pas, en programmation on appelle ça le "backtracking".
Pour une 10aines de pièce il trouve une solution en 1sec mais pour 30 pièces il lui faut une vingtaine de minutes et pour 50 pièces il faudrait plusieurs décennies !

Les informaticiens étudient les problèmes faciles à résoudre et les problèmes difficile à résoudre pour savoir si on peut trouver une meilleure solution: parfois on peut trouver une manière plus rapide, mais parfois ce on n'y parvient pas: c'est le cas de notre problème, mais aussi du sudoku et de plein d'autres problèmes.
L'étude de la difficulté à résoudre un problème s'appelle la "théorie de la complexité", et il y a encore des questions sans réponse: on ne sait toujours pas s'il existe une manière rapide de résoudre notre problème ou si cela est impossible (P=NP)

C'est très important, car on utilise le fait que certain problème sont difficiles à résoudre (comme factoriser des nombres) en cryptographie, pour protéger nos cartes bancaires et nos achats internet.

En résumé, la théorie de la complexité est l'étude du temps que certain problème mettent à être résolu par un ordinateur, certain pouvant être résolu rapidement, tandis que pour d'autre on n'arrive pas à trouver des algorithmes efficaces, même si le problème à des solutions simples à vérifier.

	
	\paragraph*{Étayages} (si les élèves patinent) :
	\begin{itemize}
		\item Placer une des pièces correctement
		\item Faire des plus petit puzzle
		\item Pour le robot: on autorise un peu plus que des instructions élémentaire
	\end{itemize}
	
	\paragraph*{Extensions} (si certains groupes d'élèves vont trop vite) :
	\begin{itemize}
        \item Il y a 6 puzzles différents
		\item Insister sur l'algorithme
	\end{itemize}
\end{document}
